Este trabajo se presentó como ponencia oral en el 12.º Encuentro de Jóvenes Investigadores
del Estado de Chihuahua, con el título «Sistema distribuido de memorias heteroasociativas
entrópicas especializadas con una arquitectura de enrutamiento semántico y recuperación
asociativa multimodal basada en la teoría de memoria transactiva».

El mismo enfoque, memorias asociativas entrópicas organizadas como una memoria
transactiva, se aplicó además a la detección de fallas en piezas industriales.
Ese trabajo está en preparación para enviarse a la revista CULCYT de la Universidad
Autónoma de Ciudad Juárez.
